\ifdefined\mainfile\else
\documentclass{article}
\usepackage{graphicx}
\usepackage{float}
\usepackage{amsmath}
\usepackage{amssymb}
\usepackage[a4paper, left=2.5cm, right=2.5cm, top=2.5cm, bottom=2.5cm]{geometry}
\usepackage[colorlinks=true,linkcolor=black,citecolor=blue,urlcolor=blue]{hyperref}
\begin{document}
\fi

\section{How we present data to the user}
\label{sec:presentation}

The user sees information about the bin in three different places, each chosen for a different context.

\subsection{The web dashboard}
\label{sec:dashboard-ui}

The primary surface is the dashboard at \url{https://dtutrash.netlify.app/}. It is intentionally a single page with four large cards -- \emph{Bags remaining}, \emph{Distance}, \emph{State} and \emph{Temperature} -- so the user can glance at it on their phone and immediately understand the bin's status. The page auto-refreshes every 5 seconds.

Two small UX details are worth calling out because they cost real effort to get right:

\begin{itemize}
    \item \textbf{ONLINE/OFFLINE pill.} The pill in the top-right corner is not based on HTTP success -- ThingSpeak will happily serve cached data when the ESP is unplugged. Instead the dashboard checks the age of the latest ESP reading: if it is older than 90\,s (three times the ESP's polling cycle), the pill flips to OFFLINE. This honestly reflects whether the bin is alive.
    \item \textbf{Filtering out command-only entries.} When the user writes a new bag count, ThingSpeak creates a new entry where \texttt{field5} is set but \texttt{field1/2/3} are \texttt{null}. If the dashboard simply showed the latest entry it would briefly display \texttt{--} for distance and bags. Instead it pulls the 10 most recent entries and walks backwards to find the latest one with \texttt{field1 != null}.
\end{itemize}

The dashboard also has a single input box and a Save button to set the bag inventory remotely. The flow is: the dashboard writes the desired count to field 5; the ESP polls field 5 within 30\,s, copies it into its internal counter, and confirms by writing the new value back to field 2 on the next push. The dashboard then sees field 2 update on its next refresh. Worst-case round-trip is approximately $30 + 16 + 5 = 51$ s, which we communicate to the user with a toast that says "Set to X bags -- ESP fetches in $<$30\,s".

\subsection{The on-device OLED}
\label{sec:oled}

A 0.96\,\textquotedbl\ SSD1306 OLED on the lid shows the current FSM state in plain English (\textsc{Ready}, \textsc{Checking}, \textsc{Closing bag}, \textsc{Replace bag}) plus the current bag inventory and the latest temperature reading. This is the channel for moments when the user is actually standing at the bin and does not want to pull out their phone.

\subsection{Push notifications}
\label{sec:notifications}

For asynchronous alerts -- "your bag is full" or "you are running out of bags" -- the system pushes through IFTTT. The user does not need any app of ours installed; they only need IFTTT, which most households already have for other automations. Two notifications are configured:

\begin{itemize}
    \item \emph{Trash bag is full -- please replace it} (triggered by \texttt{field4 = 1}).
    \item \emph{Running low on trash bags} (triggered by \texttt{field2 $<$ 3}).
\end{itemize}

This three-channel split -- live dashboard for monitoring, OLED for at-the-bin glances, notifications for time-critical events -- mirrors how the user would actually interact with the product in a real kitchen.

\ifdefined\mainfile\else
\end{document}
\fi
